% Компілювати XeLaTeX
\documentclass[11pt,a4paper]{article}
\usepackage{fontspec}
\usepackage[ukrainian]{babel}
\usepackage[margin=16mm]{geometry}
\usepackage{array,tabularx,enumitem}
\setmainfont{Arial}
\setlength{\parindent}{0pt}
\setlength{\parskip}{5pt}
\renewcommand{\arraystretch}{1.35}
\newcolumntype{Y}{>{\raggedright\arraybackslash}X}
\newcommand{\header}{%
\begin{center}\Large\textbf{Діагностична робота з STEM}\par\normalsize 5 клас\end{center}
Прізвище, ім'я \rule{0.59\linewidth}{0.4pt}\hfill Клас \rule{0.12\linewidth}{0.4pt}\par
Дата \rule{0.18\linewidth}{0.4pt}\hfill Максимум: \textbf{12 балів}\par\vspace{2pt}\hrule\vspace{4pt}}
\newcommand{\stemtask}[2]{\textbf{Завдання 1 (3 бали).} Для кожної з букв \textbf{#1} і \textbf{#2} наведи \textbf{по 2 приклади професій}, пов'язаних із відповідним складником STEM.\par
\begin{tabularx}{\linewidth}{|Y|Y|}\hline
\textbf{Буква #1:} \rule{0.35\linewidth}{0.4pt} & \textbf{Буква #2:} \rule{0.35\linewidth}{0.4pt}\\[7pt]\hline
\rule{0.4\linewidth}{0.4pt} & \rule{0.4\linewidth}{0.4pt}\\[7pt]\hline
\end{tabularx}}
\newcommand{\professions}[4]{\textbf{Завдання 2 (3 бали).} Розподіли 10 професій у три колонки. Запиши кожну професію лише один раз.\par
{\small\textit{Професії:} #1.\par}
\begin{tabularx}{\linewidth}{|Y|Y|Y|}\hline
\textbf{Людина---#2} & \textbf{Людина---#3} & \textbf{Людина---#4}\\\hline
\rule{0.27\linewidth}{0.4pt} & \rule{0.27\linewidth}{0.4pt} & \rule{0.27\linewidth}{0.4pt}\\[5pt]
\rule{0.27\linewidth}{0.4pt} & \rule{0.27\linewidth}{0.4pt} & \rule{0.27\linewidth}{0.4pt}\\[5pt]
\rule{0.27\linewidth}{0.4pt} & \rule{0.27\linewidth}{0.4pt} & \rule{0.27\linewidth}{0.4pt}\\[5pt]
\rule{0.27\linewidth}{0.4pt} & \rule{0.27\linewidth}{0.4pt} & \rule{0.27\linewidth}{0.4pt}\\[5pt]\hline
\end{tabularx}}
\newcommand{\energy}[3]{\textbf{Завдання 3 (3 бали).} У кабінеті протягом \textbf{#1 годин} працювали: \textbf{#3}. Заповни таблицю. Обчисли загальну кількість спожитої електроенергії та її вартість за тарифом \textbf{#2 грн за 1 кВт\,год}.\par
\begin{tabularx}{\linewidth}{|Y|>{\centering\arraybackslash}p{24mm}|>{\centering\arraybackslash}p{27mm}|>{\centering\arraybackslash}p{35mm}|}\hline
\textbf{Назва пристрою} & \textbf{Кількість} & \textbf{Потужність} & \textbf{Кількість спожитої електроенергії за #1 год}\\\hline
 & & & \\[12pt]\hline
 & & & \\[12pt]\hline
 & & & \\[12pt]\hline
\textbf{Разом} & & & \\\hline
\end{tabularx}\par
Вартість електроенергії: \rule{0.43\linewidth}{0.4pt} грн.}
\newcommand{\projecttask}[1]{\textbf{Завдання 4 (3 бали).} Склади короткий план проєкту \textbf{«Створення #1»}. Запиши мету, ідею, план роботи та потрібні ресурси.\par
\begin{tabularx}{\linewidth}{|Y|Y|}\hline
\textbf{Мета проєкту:} \rule{0.29\linewidth}{0.4pt} & \textbf{Ідея проєкту:} \rule{0.29\linewidth}{0.4pt}\\[10pt]\hline
\textbf{План роботи (3--4 кроки):}\par 1. \rule{0.33\linewidth}{0.4pt}\par 2. \rule{0.33\linewidth}{0.4pt}\par 3. \rule{0.33\linewidth}{0.4pt} & \textbf{Ресурси (матеріали, інструменти, люди):}\par \rule{0.38\linewidth}{0.4pt}\par \rule{0.38\linewidth}{0.4pt}\par \rule{0.38\linewidth}{0.4pt}\\[2pt]\hline
\end{tabularx}}
\begin{document}

\header
\section*{Варіант 1}
\stemtask{S}{T}
\professions{агроном, ветеринар, лісник, садівник, художник, дизайнер, фотограф, актор, програміст, механік}{природа}{художній образ}{техніка}
\energy{8}{4,32}{6 ноутбуків потужністю 60 Вт кожен, 1 проєктор потужністю 200 Вт і 4 настільні лампи потужністю 10 Вт кожна}
\projecttask{мультфільму}

\newpage\header
\section*{Варіант 2}
\stemtask{E}{M}
\professions{еколог, бджоляр, зоотехнік, гідролог, музикант, ілюстратор, модельєр, режисер, електрик, інженер-конструктор}{природа}{художній образ}{техніка}
\energy{8}{4,32}{5 комп'ютерів потужністю 120 Вт кожен, 1 інтерактивна дошка потужністю 180 Вт і 1 принтер потужністю 100 Вт}
\projecttask{самоката}

\newpage\header
\section*{Варіант 3}
\stemtask{S}{E}
\professions{фермер, квітникар, рибовод, геолог, скульптор, аніматор, артист, перукар, оператор верстатів, фахівець з ремонту телефонів}{природа}{художній образ}{техніка}
\energy{8}{4,32}{8 планшетів потужністю 15 Вт кожен, 1 телевізор потужністю 100 Вт і 5 зарядних пристроїв потужністю 12 Вт кожен}
\projecttask{машини}

\newpage\header
\section*{Варіант 4}
\stemtask{T}{M}
\professions{кінолог, ґрунтознавець, садівник, пасічник, хореограф, дизайнер одягу, художник-мультиплікатор, ведучий, майстер з ремонту велосипедів, робототехнік}{природа}{художній образ}{техніка}
\energy{8}{4,32}{7 моніторів потужністю 25 Вт кожен, 7 ноутбуків потужністю 60 Вт кожен і 2 колонки потужністю 20 Вт кожна}
\projecttask{парку атракціонів}

\end{document}
